\documentclass[11pt]{article}

\usepackage[margin=0.9in]{geometry}
\usepackage[T1]{fontenc}
\usepackage[utf8]{inputenc}
\usepackage{booktabs}
\usepackage{enumitem}
\usepackage{graphicx}
\usepackage{hyperref}
\usepackage{xcolor}

\hypersetup{
  colorlinks=true,
  linkcolor=blue!55!black,
  urlcolor=blue!55!black,
  citecolor=blue!55!black
}

\setlist[itemize]{leftmargin=1.4em}
\setlist[enumerate]{leftmargin=1.6em}
\setlength{\emergencystretch}{3em}

\newcommand{\pzr}{Predictive Zonotope Reduction}
\newcommand{\mpcg}{\textsc{MPC-G}}
\newcommand{\mpcw}{\textsc{MPC-W}}

\title{Predictive Zonotope Reduction: Deep-Research Brief}
\author{Project notes synthesized from \texttt{science/SCIENCE.md} and \texttt{results/tacas-main}}
\date{May 2026}

\begin{document}
\maketitle

\begin{abstract}
This brief frames \pzr{} as a research problem for follow-up literature review
and hypothesis refinement. The project starts from stream-based runtime
monitoring under calibration and measurement uncertainty. Affine arithmetic
tracks such uncertainty symbolically, but fresh measurement noise introduces
new slack variables at each step, so the exact monitor state grows without
bound. Existing zonotope reductions solve the bounded-memory problem by
over-approximating the state, but they usually optimize current-set geometry
rather than future monitor verdict quality. The project idea is to treat
zonotope reduction as a certified abstraction-control action: an untrusted
static, predictive, or learned policy chooses among certified reducers, while
the reducer certificate alone preserves soundness and bounded memory.
\end{abstract}

\section{Purpose of This Document}

This document is designed as input for a deep-research pass. It does not try to
be a paper draft. Instead, it records the project's motivation, technical
contract, implementation shape, current empirical status, and the literature
questions that need better grounding before the work is written for TACAS or a
similar venue.

The companion advisor memo generated under
\path{results/tacas-main/advisor-report/} contains the numerical tables for the
latest evaluation. This brief focuses on why the problem is interesting, how
the current prototype maps to prior work, and where deeper research should look
for novelty, threats, and stronger comparisons.

\section{Problem Motivation}

Stream-based runtime monitors transform input streams into derived output
streams and trigger verdicts. In safety-critical settings, sensor measurements
are rarely exact: some errors are persistent calibration biases, while others
are independent measurement noise. The paper ``Cutting Corners on Uncertainty''
shows that affine arithmetic can track these uncertainties in monitor state,
and that the affine coefficients form a zonotope
\[
  Z = c + G[-1,1]^m.
\]
Here, the center \(c\) stores nominal stream values and each generator column
in \(G\) stores one symbolic uncertainty direction. Persistent calibration
error can reuse a symbolic variable, but independent measurement noise creates
a fresh generator at each step. Therefore, an exact monitor state grows with
trace length.

The bounded-memory problem is not optional. Runtime monitors are intended to
run continuously beside a system. A trace-length-dependent memory footprint is
at odds with the usual runtime-assurance expectation that monitor resource
requirements are statically bounded. The baseline solution is to reduce or
unify generators online, producing a zonotope that encloses the exact affine
state while keeping at most \(K\) generators.

\section{Core Idea}

The project's central move is to stop treating reduction as a fixed local
operator and instead treat it as an abstraction-control decision. At each
over-budget state, a policy chooses an action from a finite set of certified
reducers:
\[
  a_t = \pi(H_t, Z_t, \widehat{x}_{t:t+h}),
  \qquad
  Z'_t = \rho_{a_t,K}(Z_t).
\]
The policy can use the current zonotope, trigger geometry, metadata about
calibration and measurement generators, and a short predicted future trace. In
the predictive variants, it evaluates candidate first reductions by rolling the
monitor forward over predicted inputs and scoring future trigger imprecision.
Only the first action is applied, then the monitor replans at the next real
overflow.

This is model-predictive control in structure, but the control input is not a
physical actuator. It is a compression choice. The objective is not tracking a
desired plant trajectory; it is preserving future monitor precision under a
fixed generator budget.

\section{Trusted Boundary}

The trusted boundary is deliberately small:
\[
  Z \subseteq \rho_{a,K}(Z),
  \qquad
  \mathrm{gen}(\rho_{a,K}(Z)) \le K.
\]
If every candidate reducer returns a certificate for this contract, then any
policy selecting among those candidates preserves soundness. The selector can
be heuristic, approximate, learned, adversarially bad for precision, or based
on an inaccurate predictor; it still cannot make the monitor unsound unless it
is allowed to bypass the certified reducer contract.

This is the core theorem shape to preserve in the eventual paper:

\begin{quote}
Assume the unreduced symbolic monitor transition soundly encloses all concrete
states represented by its input zonotope. If every applied reduction returns a
sound bounded over-approximation, then the reduced execution over-approximates
the unreduced execution independently of the policy used to choose reducers.
\end{quote}

The implementation also tracks generator metadata. Calibration, measurement,
synthetic, and unknown generator columns are tagged. A monitor can declare
generator requirements; \texttt{ProtectedReducer} copies required generators
exactly into the reduced state and only compresses the residual. If the required
columns alone exceed the budget, the reducer fails rather than silently
dropping semantically required uncertainty.

\section{Implementation Map}

The Python prototype is organized around four boundaries:

\begin{itemize}
\item \textbf{Core zonotopes and certificates.} \texttt{pzr.core} stores
zonotope values, generator metadata, and reduction certificates.
\item \textbf{Reducers.} \texttt{pzr.reduction} implements static and
monitor-aware reducers including box, Girard, Combastel, MethA, Scott, PCA,
adaptive, scored-keep, identity, target-budget, and protected wrappers.
\item \textbf{Monitor adapters.} \texttt{pzr.monitoring} exposes a black-box
boundary: the controller can step, clone, inspect trigger metadata, and replace
the zonotope component, but it does not depend on monitor equations.
\item \textbf{Control policies.} \texttt{pzr.control} contains static,
one-action MPC, sequence MPC, and rollout MPC selectors.
\end{itemize}

The current benchmark scenarios are:

\begin{itemize}
\item \texttt{robot}: the harder omnidirectional robot monitor, tracking
filtered acceleration, velocity, distance increment, and x/y position.
\item \texttt{robot\_simple}: the two-axis motivating robot with persistent
x/y calibration generators and fresh measurement generators.
\item \texttt{thermostat}: a non-robot monitor with temperature, filtered
temperature, heating/cooling effort, comfort deviation, and safety triggers.
\end{itemize}

There are two naming layers to keep straight. Code-level method IDs include
\path{mpc_focused_fixed_girard} and
\path{mpc_wide_fixed_girard}. Latest TACAS-main aggregate artifacts expose
the corresponding report-facing names \texttt{mpc\_rollout\_girard} and
\texttt{mpc\_rollout\_wide}. In this document, \mpcg{} refers to the focused
fixed-Girard rollout and \mpcw{} to the wide first-action rollout.

\begin{table}[htbp]
\centering
\caption{Current method roles.}
\begin{tabular}{lll}
\toprule
Method family & Role & Soundness dependency \\
\midrule
Static reducers & Baseline local compression & Reducer certificate \\
\mpcg{} & Focused predictive first-action search & Protected certified reducer \\
\mpcw{} & Broad first-action ablation & Protected certified reducer \\
Learned distilled & Cheap selector approximation & Certified fallback/reducer set \\
Reference & Unbounded precision reference & Not a bounded competitor \\
\bottomrule
\end{tabular}
\end{table}

\section{Background Work to Investigate}

\subsection{Stream-Based Runtime Monitoring}

The stream-monitoring lineage is the semantic home for this project. Lola and
RTLola model monitors as stream transformations, with trigger streams used for
verdicts. Recent RTLola work emphasizes automatic monitor construction, static
memory bounds, runtime safety, and use in cyber-physical settings. The
uncertainty paper extends this setting by treating sensor values as uncertain
quantities rather than exact inputs.

Deep-research questions:

\begin{itemize}
\item How have stream-monitoring languages formalized bounded-memory
requirements, and can our generator budget be presented as a compatible
resource bound?
\item Are there prior runtime-monitoring papers that optimize monitor
abstractions dynamically, rather than statically compiling a fixed abstraction?
\item What terminology will be legible to TACAS reviewers: approximation
control, abstraction control, online abstraction refinement, or monitor-aware
order reduction?
\end{itemize}

\subsection{Zonotope Order Reduction}

Classical zonotope order reduction aims to enclose a high-order zonotope by a
lower-order one with acceptable over-approximation. Girard-style,
Combastel-style, PCA-based, and Scott-style reductions come from this
literature. Yang and Scott compare reduction techniques by computational cost
and overestimation error; Kopetzki, Schuermann, and Althoff survey known
methods and introduce PCA, clustering, and optimization-based variants.

The gap for this project is semantic and temporal. These methods generally
optimize current set geometry. A runtime monitor cares about future trigger
verdicts, persistent calibration variables, and the way old approximation error
flows through future stream equations.

Deep-research questions:

\begin{itemize}
\item Which order-reduction metrics in the control literature map cleanly to
our interval-hull MSE, trigger width, and false-alarm metrics?
\item Is there prior work on task-aware or property-aware zonotope reduction,
especially where the downstream task is verification or monitoring?
\item Can protected generator preservation be related to constrained zonotopes,
structured uncertainty, or parameter-varying set-membership estimation?
\end{itemize}

\subsection{Set-Membership Estimation and Robust Control}

Zonotopes are common in set-membership estimation, reachability analysis,
fault diagnosis, and robust control because they are compact and closed under
linear maps and Minkowski sums. This literature gives the mathematical
vocabulary for guaranteed enclosures and bounded-complexity propagation.

The connection to this project is that monitor state behaves like a symbolic
state estimator over derived quantities. However, our reduction objective is
not plant-state estimation error in the usual sense. It is preservation of
monitor verdict precision under stream equations and trigger predicates.

\subsection{MPC and Robust/Tubed Prediction}

\pzr{} borrows MPC's receding-horizon pattern: optimize over a finite predicted
horizon, apply the first action, then replan. Robust MPC and tube MPC are
natural analogies for future work if prediction uncertainty becomes important.
At present, the latest online-versus-oracle artifacts do not show a large
consistent oracle advantage, so robust prediction is better framed as a future
research direction rather than a current empirical claim.

Deep-research questions:

\begin{itemize}
\item Is there existing theory for MPC over abstractions or compression
operators rather than plant inputs?
\item Can rollout MPC with a fixed tail policy be described using standard
approximate dynamic programming or rollout-control terminology?
\item Would a robust predictor improve selector stability, or is the main
bottleneck currently the reducer action set and scoring objective?
\end{itemize}

\section{Current Empirical Status}

The latest artifact inspected is \path{results/tacas-main}. Its manifest
records a paper-profile run with length 300, budget 8, horizon 6, 50 seeds,
both online and oracle predictor modes, and method set
\path{paper_plus_wide}. It covers \texttt{robot}, \texttt{robot\_simple}, and
\texttt{thermostat}. The aggregate diagnostics are
clean: zero budget violations, zero unsound certificates, zero reduction
failures, and zero explicit no-op selections.

The most defensible empirical story is:

\begin{itemize}
\item On the hard robot monitor, \mpcg{} improves precision over strong static
baselines at the same budget. In online mode at budget 8, mean trigger width is
225.1 for \mpcg{} versus 230.1 for Girard; mean interval-hull MSE is 268.4 for
\mpcg{} versus 352.6 for Girard.
\item In the robot budget sweep, \mpcg{} beats Girard from budgets 8 through
20. This is the cleanest current figure-level evidence for predictive
reduction.
\item \mpcw{} is not the headline method. It searches a broader first-action
space and evaluates many more sequences, but on the hard robot it is less
precise than \mpcg{} and Girard.
\item The learned policy reaches about 90.6\% validation accuracy and 98.8\%
top-3 accuracy in the latest artifact, but the training labels are heavily
Girard-dominated. It should be framed as promising infrastructure, not as the
main scientific result.
\item Robot-simple and thermostat are useful sanity checks, but several
certified reducers saturate them. They should not be used to claim broad
dominance.
\end{itemize}

\section{Likely Contribution Claim}

The safest paper claim is:

\begin{quote}
Bounded-memory monitoring under sensor uncertainty benefits from policy-guided
certified zonotope reduction. The policy can be predictive or learned, while
soundness remains policy-independent because only certified reducers may mutate
the monitor state. A focused receding-horizon selector improves precision on a
hard robot monitor compared with strong static order-reduction baselines.
\end{quote}

This claim avoids overcommitting on three weak points in the current evidence:
\mpcw{} is not dominant, learned selection is not yet a strong independent
contribution, and oracle prediction is not currently essential.

\section{Deep-Research Agenda}

\begin{enumerate}
\item \textbf{Positioning.} Find the closest prior work on task-aware set
approximation, monitor-aware abstraction, online abstraction control, and
property-preserving order reduction.
\item \textbf{Formalization.} Express the policy-independent soundness theorem
using a clean monitor transition relation, a reducer action contract, and a
bounded-memory invariant.
\item \textbf{Action-set design.} Explain why the focused rollout action set
works better than broad first-action search in the current artifacts. Search
for theory or empirical precedent about action-space restriction in rollout
control.
\item \textbf{Learning.} Decide whether the learned selector is an evaluation
side story or a contribution. A stronger version likely needs balanced labels,
focused-rollout labels, or a learned scorer rather than a direct classifier.
\item \textbf{Prediction.} Determine whether robust or tube-aware prediction is
scientifically worth adding. The current oracle gap suggests caution.
\item \textbf{Benchmarks.} Add or tune a non-robot benchmark where static
reducers do not saturate, so the evaluation is not robot-specific.
\end{enumerate}

\section{Search Terms for Deep Research}

\begin{itemize}
\item stream-based runtime monitoring uncertainty affine arithmetic zonotope
\item RTLola bounded memory runtime monitoring sensor uncertainty
\item zonotope order reduction property aware task aware
\item generator reduction monitor precision false alarms
\item set-membership estimation zonotope reduction structured uncertainty
\item model predictive control abstraction refinement
\item rollout policy fixed tail action set approximate dynamic programming
\item robust MPC tube prediction set-valued future uncertainty
\item runtime verification quantitative uncertainty monitor abstraction
\end{itemize}

\section{Claim Boundaries}

Use the following boundaries when turning this into paper prose:

\begin{itemize}
\item Do claim policy-independent soundness under certified reducers.
\item Do claim focused rollout improves hard-robot precision in current
artifacts.
\item Do not claim \mpcw{} is always better; it is a negative/diagnostic
ablation.
\item Do not claim learned selection is the main empirical advance yet.
\item Do not claim oracle prediction is necessary; current evidence is mixed
and small.
\item Do not treat the unbounded reference as a bounded-memory competitor.
\end{itemize}

\section{Selected Source Anchors}

\begin{thebibliography}{9}

\bibitem{cutting-corners}
Bernd Finkbeiner, Martin Fraenzle, Florian Kohn, and Paul Kroeger.
\newblock Cutting Corners on Uncertainty: Zonotope Abstractions for
Stream-based Runtime Monitoring.
\newblock arXiv:2601.11358, 2026.
\newblock \url{https://arxiv.org/abs/2601.11358}

\bibitem{yang-scott}
Xuejiao Yang and Joseph K. Scott.
\newblock A comparison of zonotope order reduction techniques.
\newblock \emph{Automatica}, 95:378--384, 2018.
\newblock \url{https://doi.org/10.1016/j.automatica.2018.06.006}

\bibitem{kopetzki}
Anna-Kathrin Kopetzki, Bastian Schuermann, and Matthias Althoff.
\newblock Methods for Order Reduction of Zonotopes.
\newblock IEEE CDC, 2017.
\newblock \url{https://doi.org/10.1109/CDC.2017.8264508}

\bibitem{combastel}
Christophe Combastel.
\newblock A state bounding observer based on zonotopes.
\newblock European Control Conference, 2003.
\newblock \url{https://skoge.folk.ntnu.no/prost/proceedings/ecc03/pdfs/437.pdf}

\bibitem{rtlola}
Stream-based monitoring with RTLola.
\newblock \emph{Science of Computer Programming}, 253, 2026.
\newblock \url{https://doi.org/10.1016/j.scico.2026.103495}

\bibitem{tacas-artifacts}
ETAPS/TACAS 2026 conference and artifact information.
\newblock \url{https://etaps.org/2026/conferences/tacas/}

\end{thebibliography}

\end{document}
